\backmatter
\silentchapter{Etterord}{etterord}
\begin{widebody}
Empirien er tufta på eit korpus av om lag 2 300 europeiske stolar (1280–2024). Datasettet tillet kvantitativ testing av formhistoriske hypotesar. Analysen av meshgeometri har synt at sju mesh-avleia trekk aukar den prediktive krafta for stilperiode med ein faktor på over fire samanlikna med katalogdimensjonar.

Overgangen frå prosa til proposisjonsform var ein logisk nødvendigheit. Desimalnummereringa angjev den logiske vekta, og superscript (d, a, t, o, i) indikerer logisk status. Til skilnad frå Wittgensteins original, er kvar proposisjon knytt til eit eksplisitt falsifiseringsvilkår; traktaten er konstruert for å kunne falsifiserast. Dei empiriske testane stadfestar hovudproposisjonane: formrommet er ikkje-uniformt busett; stilperiode er ein effektiv proxy for aggregerte seleksjonstrykk; og landskapet er i uopphøyrleg endring. Funna er robuste på tvers av ulike geometriske representasjonar.

Observert latens i stålillustrasjonen viser at formhistoria er langsamare enn materialhistoria; eit nytt substrat arvar formspråket til det substratet det erstattar, heilt til eigne affordansar pressar det ut i nye regionar.

Traktaten skildrar krefter og posisjonar, ikkje verdiar. Ho kan forklare avstanden i formrommet, men ho kan ikkje felle estetiske domar. Estetikk er transcendent til formsystemet, på same vis som etikk er transcendent til logikken (Wittgenstein). Agenten navigerer; modellen registrerer berre den resulterande posisjonen. Proposisjonane skal erkjennast som ein stige. Dei skal ikkje erstatte handverket, men gjere den intuitive navigasjonen eksplisitt ved å vise agenten kva han allereie gjer.

At teksten kan falsifiserast, er garantien for hennar gyldigheit. Om eitt einaste seleksjonstrykk forklarer all formvariasjon, kollapsar den logiske kjeden. Traktaten er stigen som skal kastast.

Takk til Jan Petter Neverdahl og Hermann Stange for kommentarar under arbeidet.

\end{widebody}
\silentchapter{Notatar}{notatar}
\begin{widebody}
\theendnotes
\end{widebody}

\silentchapter{Referansar}{referansar}
\begin{widebody}
\refentry{Sullivan, L. H. (1896). The tall office building artistically considered. Lippincott's Magazine, 57, 403–409.}
\refentry{Thompson, D'A. W. (1917). On growth and form. Cambridge University Press.}
\refentry{Wright, S. (1932). The roles of mutation, inbreeding, crossbreeding, and selection in evolution. Proc. Sixth Int. Congress of Genetics, 1, 356–366.}
\refentry{Rosenblueth, A., Wiener, N. \& Bigelow, J. (1943). Behavior, purpose and teleology. Philosophy of Science, 10(1), 18–24.}
\refentry{Shannon, C. E. (1948). A mathematical theory of communication. Bell System Technical Journal, 27, 379–423.}
\refentry{Wiener, N. (1948). Cybernetics. MIT Press.}
\refentry{Turing, A. M. (1950). Computing machinery and intelligence. Mind, 59(236), 433–460.}
\refentry{Waddington, C. H. (1957). The strategy of the genes. George Allen \& Unwin.}
\refentry{Kubler, G. (1962). The shape of time. Yale University Press.}
\refentry{Kuhn, T. S. (1962). The structure of scientific revolutions. University of Chicago Press.}
\refentry{Raup, D. M. (1966). Geometric analysis of shell coiling. Journal of Paleontology, 40(5), 1178–1190.}
\refentry{Eldredge, N. \& Gould, S. J. (1972). Punctuated equilibria. I T. J. M. Schopf (Red.), Models in paleobiology, 82–115.}
\refentry{Stiny, G. \& Gips, J. (1972). Shape grammars and the generative specification of painting and sculpture. Information Processing 71, 1460–1465.}
\refentry{Gibson, J. J. (1979). The ecological approach to visual perception. Houghton Mifflin.}
\refentry{Stiny, G. (1980). Introduction to shape and shape grammars. Environment and Planning B, 7(3), 343–351.}
\refentry{Stiny, G. (1991). The algebras of design. Research in Engineering Design, 2, 171–181.}
\refentry{Kauffman, S. A. (1993). The origins of order. Oxford University Press.}
\refentry{Arthur, W. B. (1994). Increasing returns and path dependence in the economy. University of Michigan Press.}
\refentry{Michl, J. (1995). Form follows WHAT? The modernist notion of function as a carte blanche.}
\refentry{Odling-Smee, F. J., Laland, K. N. \& Feldman, M. W. (2003). Niche construction. Princeton University Press.}
\refentry{Hatchuel, A. \& Weil, B. (2003). A new approach of innovative design: An introduction to C-K theory. Proc. ICED 2003.}
\refentry{Hatchuel, A. \& Weil, B. (2009). C-K design theory: An advanced formulation. Research in Engineering Design, 19(4), 181–192.}
\refentry{Stiny, G. (2006). Shape: Talking about seeing and doing. MIT Press.}
\refentry{Mitteroecker, P. \& Huttegger, S. M. (2009). The concept of morphospaces in evolutionary and developmental biology. Biological Theory, 4(1), 54–67.}
\refentry{Doctor, T., Wakefield, E., Pavlic, T. P. \& Levin, M. (2024). Biology, Buddhism, and AI: Care as the driver of intelligence. Entropy, 26(4), 352.}
\refentry{Fields, C. \& Levin, M. (2022). Competency in navigating arbitrary spaces. Entropy, 24(6), 819.}
\refentry{Levin, M. (2022). Technological approach to mind everywhere. Frontiers in Systems Neuroscience, 16, 768201.}
\refentry{Levin, M. (2025). Living things are not machines (also, they totally are). Noema Magazine.}
\end{widebody}

